\subsection{R}

\begin{frame}
	R language and environment for statistical computing.
\end{frame}

\begin{frame}
	\frametitle{Make script take parameters}
\end{frame}

\begin{frame}
	\frametitle{Iteration}
	\textit{Of course someone has to write loops. It doesn't have to be
		you}
	Jennifer Bryan
\end{frame}

\begin{frame}
	\frametitle{Iteration}
	Recomended reading:
	\begin{itemize}
		\item Iteration, chapter 26~\cite{wickham2023a}.
		\item Functionals, chapter 9~\cite{wickham2019}.
	\end{itemize}
\end{frame}

\begin{frame}
	\frametitle{Parallelize code}
\end{frame}

\begin{frame}
	\frametitle{Resources}
	\begin{itemize}
		\item Tutorials: \href{https://www.w3schools.com/r/default.asp}
		      {R Tutorial} by w3schools.
		\item Books:
		      \begin{itemize}
			      \item An introduction to R~\cite{venables2024},
			      \item Advanced R~\cite{wickham2015a}.
		      \end{itemize}
	\end{itemize}
\end{frame}

\begin{frame}
	\frametitle{Parallelize code example 1}
	\begin{itemize}
		\item Use Virtual Rastes (VRT).
		\item Split input data while avoiding breaking spatial and
		      temporal dependencies.
		\item Use functions.
		\item Functions should have no collateral effects.
		\item Use lists and lapply.
		\item Convert lapply to parallel::mclapply.
	\end{itemize}
\end{frame}

\begin{frame}
	\frametitle{R functions}
\end{frame}
